\begin{lemma}[2.2.2]
\phantomsection
\label{II.2.2.2-ko}
$d$, $e$를 $>0$인 두 정수, $f\in S_d$, $g\in S_e$라 하자. 다음과 같은
환의 표준 동형이 존재한다.
\[
  S_{(fg)}\xrightarrow{\sim}(S_{(f)})_{g^d/f^e}~;
\]
이 두 환을 표준적으로 동일시하면, 다음과 같은 가군의 표준 동형이
존재한다.
\[
  M_{(fg)}\xrightarrow{\sim}(M_{(f)})_{g^d/f^e}.
\]
\end{lemma}

\begin{proof}
실제로 $fg$는 $f^e g^d$의 약수이고, 후자는 $(fg)^{de}$의 약수이므로
등급환 $S_{fg}$와 $S_{f^e g^d}$는 표준적으로 동일시된다. 다른 한편
$S_{f^e g^d}$는 $(S_{f^e})_{g^d/1}$와도 동일시되고
(\hyperref[0.1.4.6-ko]{0, I.4.6}), $f^e/1$은 $S_{f^e}$에서 가역이므로
$S_{f^e g^d}$는 $(S_{f^e})_{g^d/f^e}$와도 동일시된다. 그런데
$g^d/f^e$는 차수가 $0$인 $S_{f^e}$의 원소이다. 따라서
$(S_{f^e})_{g^d/f^e}$에서 차수가 $0$인 원소들로 이루어진 부분환은
$(S_{(f^e)})_{g^d/f^e}$임을 곧바로 얻는다. 또 명백히
$S_{(f^e)}=S_{(f)}$이므로, 이는 명제의 첫 번째 부분을 증명한다.
두 번째 부분도 같은 방식으로 확립된다.
\end{proof}

\begin{env}[2.2.3]
\phantomsection
\label{II.2.2.3-ko}
(\hyperref[II.2.2.2-ko]{2.2.2})의 가정 아래에서, 표준 준동형
$S_f\to S_{fg}$ (\hyperref[0.1.4.1-ko]{0, I.4.1})은 $x/f^k$를
$g^k x/(fg)^k$로 보내며 차수가 $0$임은 명백하다. 따라서
제한함으로써
\emph{표준 준동형 $S_{(f)}\to S_{(fg)}$}을 얻으며, 다음 도식은
가환한다.
\[
  \xymatrix{
    & S_{(f)}\ar[dl]\ar[dr]\\
    S_{(fg)}\ar[rr]^-{\sim} && (S_{(f)})_{g^d/f^e}
  }
\]
마찬가지로 표준 준동형 $M_{(f)}\to M_{(fg)}$을 정의한다.
\end{env}

\begin{lemma}[2.2.4]
\phantomsection
\label{II.2.2.4-ko}
$f$, $g$가 $S_+$의 두 동차 원소이면, 환 $S_{(fg)}$는 $S_{(f)}$와
$S_{(g)}$의 표준 상들의 합집합에 의해 생성된다.
\end{lemma}

\begin{proof}
(\hyperref[II.2.2.2-ko]{2.2.2})에 따라
$1/(g^d/f^e)=f^{d+e}/(fg)^d$가 $S_{(g)}$의 $S_{(fg)}$ 안에서의 표준
상에 속함을 보이면 충분하며, 이는 정의상 명백하다.
\end{proof}

\begin{proposition}[2.2.5]
\phantomsection
\label{II.2.2.5-ko}
$d$를 $>0$인 정수, $f\in S_d$라 하자. 그러면 다음과 같은 환의
표준 동형이 존재한다.
$S_{(f)}\xrightarrow{\sim}S^{(d)}/(f-1)S^{(d)}$~; 이 동형으로 두 환을
표준적으로 동일시하면, 다음과 같은 가군의 표준 동형이 존재한다.
$M_{(f)}\xrightarrow{\sim}M^{(d)}/(f-1)M^{(d)}$.
\end{proposition}

\begin{proof}
이 동형들 중 첫 번째 것은 $x/f^n$($x\in S_{nd}$)에 원소
$\overline{x}$, 곧 $x$의 $(f-1)S^{(d)}$에 대한 잉여류를 대응시켜 정의한다.
이 사상은 잘 정의된다. 실제로 모든 양의 정수 지수에 대하여 합동식
$f^h x\equiv x\pmod{(f-1)S^{(d)}}$이 모든 $x\in S^{(d)}$에 대해
성립한다. 따라서 $f^h x=0$인 $h>0$가 있으면
\oldpage[II]{25}
$\overline{x}=0$이다. 한편 $x\in S_{nd}$이고 $x=(f-1)y$라 하자. 여기서
$y=y_{hd}+y_{(h+1)d}+\cdots+y_{kd}$이고 $y_{jd}\in S_{jd}$이며
$y_{hd}\neq 0$이다. 그러면 반드시 $h=n$이고 $x=-y_{hd}$이며, 또한
관계 $y_{(j+1)d}=fy_{jd}$가 $h\leq j\leq k-1$에 대해 성립하고
$fy_{kd}=0$이다.
이로부터 마침내 $f^{k-n+1}x=0$을 얻는다. 따라서 $x\in S_{nd}$인 원소가
나타내는 $(f-1)S^{(d)}$에 대한 각 잉여류 $\overline{x}$에 $S_{(f)}$의 원소
$x/f^n$을 대응시킬 수 있다.
앞의 고찰에 따라 이 사상은 잘 정의된다. 이렇게 정의된 두 사상이
환 준동형이며 서로 역임을 곧바로 알 수 있다. $M$에 대해서도 완전히
같은 방식으로 한다.
\end{proof}

\begin{corollary}[2.2.6]
\phantomsection
\label{II.2.2.6-ko}
$S$가 뇌터이면, 차수가 $>0$인 임의의 동차 원소 $f$에 대하여
$S_{(f)}$도 뇌터이다.
\end{corollary}

\begin{proof}
이는 (\hyperref[II.2.1.7-ko]{2.1.7})과
(\hyperref[II.2.2.5-ko]{2.2.5})에서 곧바로 따른다.
\end{proof}
